\chapter{CP/M: the Z80 Second Processor}\label{cha:cpm}

In 1984 Acorn sold a Z80 Second Processor for the BBC Micro; its purpose
was to run CP/M, the operating system that owned business computing before
the IBM PC. It sat on the Tube. So does ours. \verb|CPM| at the shell boots
CP/M~2.2 on a Z80 co-processor (RunCPM, on the host), and the \verb|A0>|
prompt that defined a decade appears. The boot banner cheerfully measures
the Z80 in \emph{gigahertz} --- read that number knowingly: it is the
emulated Z80's speed \emph{on whatever host you run the emulator on} (a
laptop, in this book's screenshots), so it varies machine to machine.
For scale, the real thing shipped at 4\,MHz.

\begin{type}
CPM
\end{type}

\screen{shots/cpm.png}{CP/M 2.2 over the Tube: the CCP's \texttt{DIR} and
\texttt{TYPE}, on drive A.}

\section{Drives are folders}
Drives \verb|A:| to \verb|P:| are the host folders \pth{fs/CPM/A}
through \pth{fs/CPM/P}; CP/M's user areas 0--15 are numbered subfolders.
Drop any \texttt{.COM} file into \pth{fs/CPM/A/0} and run it by name.
\verb|DIR|, \verb|TYPE|, \verb|ERA|, \verb|REN|, \verb|SAVE| and
\verb|USER| are built into the CCP; \verb|EXIT| hands the console back to
the K/OS shell.

Three of the drives are shortcuts rather than folders of their own:
\verb|K:| is the machine's own filesystem, so CP/M and K/OS can hand files
to each other in both directions; \verb|P:| is Turbo Pascal and \verb|D:|
is WordStar, each pointing straight at the user area the program lives in.

And one drive letter is spoken for: \verb|N:| is the network, reserved
and empty. It is one letter and not two because on this machine the
scheme lives in the \emph{name} rather than in the device --- FujiNet
and Meatloaf are two URL schemes over one namespace, not two kinds of
place --- so \verb|N:| will be a mount point showing whatever the
shell's \verb|CD tnfs://| has the machine looking at. The plumbing is
not written; the letter is held so that nothing else takes it.

\section{Starting a program from the shell}
\verb|CPM| on its own gives you the \verb|A0>| prompt. \verb|CPM command|
runs a command at boot instead, through the CCP's \texttt{AUTOEXEC.TXT}.
One line is all the CCP reads --- but a name with no \texttt{.COM} to
match runs the \texttt{.SUB} of that name, and that may be as long as you
like. \pth{/CPM/A/0/K-TURBO.SUB} is the pattern:

\begin{type}
P:
TURBO
EXIT
\end{type}

Ending in \verb|EXIT| is what makes the round trip whole: quitting a CP/M
program only drops you to the CCP, and the \verb|EXIT| carries you the
rest of the way back to the K/OS prompt. With that in place
\verb|ALIAS TURBO CPM K-TURBO| makes \verb|TURBO| a word you can type at
the shell; \pth{/SYSTEM/ETC/STARTUP.SAMPLE} defines that one, \verb|WS|
and \verb|MBASIC|.

\begin{aside}
\textbf{Why not \texttt{USER} in a submit.} The obvious script is
\verb|H:|, \verb|USER 3|, \verb|TURBO| --- and it dies on the second line.
The submit in progress is a file, \texttt{\$\$\$.SUB}, and it lives on
A: user~0; \verb|USER 3| walks away from it and it cannot be read any
more. That is CP/M's own behaviour and not RunCPM's doing. Pointing a
drive at the user area, as \verb|P:| and \verb|D:| do, keeps \verb|USER|
out of it.
\end{aside}

\section{Typing a CP/M program's name directly}
The shell can be told to treat an unknown word as a CP/M program: F7
\(\rightarrow\) Shell \(\rightarrow\) \emph{CP/M .COM by name}. With it on,
\verb|STAT| at the \verb|/]| prompt starts CP/M, runs \texttt{STAT.COM}
from A: and comes back.

It is off to begin with, and deliberately so --- \verb|D| typed for
\verb|DIR| should not start a Z80 program. It changes only the guess: a
command naming \verb|CPM| outright, an alias among them, works either way.
There is nothing in a \texttt{.prg} to tell the two apart, incidentally ---
its header is a load address and a run address, and a \texttt{.COM} begins
with Z80 code that would read as a perfectly plausible one --- so the
extension is what decides.

\section{Software}
The RunCPM project ships a complete system disk
(\texttt{DISK/A0.zip} in its repository): Digital Research's own
\texttt{ASM}, \texttt{MAC}, \texttt{DDT}, \texttt{ZSID}, \texttt{STAT},
\texttt{PIP} and \texttt{ED}, the \texttt{SUBMIT} batch system, Microsoft
BASIC, a Z80 assembler with its manual, XMODEM, and the source code of the
BDOS and CCP --- buildable on the machine itself. One
\texttt{unzip A0.zip -d fs/CPM/} installs the lot. Beyond that lies the
whole surviving CP/M software world: WordStar, Turbo Pascal, dBase~II,
Zork, one archive away.

\begin{aside}
\textbf{Edges, honestly:} line-mode programs (assemblers, compilers,
MBASIC, adventures) work today, and so does the full-screen software:
the console is JIM, a VT100 in hardware (Chapter~\ref{cha:tube}).
WordStar~4 on \texttt{E:} user~3 comes installed for ``ANSI Standard''
at 79$\times$29 --- \verb|WS| and you are editing; Turbo Pascal~3 on
\texttt{H:} user~3 needs nothing either. A program that asks its
terminal type at install time wants ``VT100'' or ``ANSI''.

The four arrow keys work in CP/M software: while the Tube is running
CP/M the machine sends them down as the WordStar keys --- Ctrl-E, X, S
and D for up, down, left and right --- which is what the programs of 1984
read, WordStar and Turbo Pascal (Chapter~\ref{cha:pascal}) among them.
The rest of the navigation keys do not: Home, End, PgUp and PgDn are
two-key \texttt{\textasciicircum Q} sequences in WordStar, and there is
no single key the machine could honestly send for them, so type those as
the program's manual says.

\screeninline{shots/wordstar.png}{WordStar~4 editing its own \texttt{READ.ME}:
the edit menu, the ruler, the text --- 79$\times$29 on JIM.}
\end{aside}
